\section{Introduction}\label{sec:intro}

At first glance, automated unit test generation using large language models seems nearly
solved. For example, GPT-4o achieves 98.65\% line coverage on
TestEval~\cite{wang2025testeval}, and GPT-4-class model panels report similar numbers on
HumanEval-derived suites~\cite{kumar2025empirical}. The functions behind those numbers,
however, are short, self-contained, LeetCode-style problems. A test that performs well
under these conditions has never had to deal with a real import graph, navigate a class
hierarchy, or handle overloaded method signatures, and high line coverage alone does not
guarantee that the test will actually catch faults. When we instead sampled 120 functions
from ten real projects --- from Apache Commons, Gson and Joda-Time to Flask and Requests
--- we found that most of the test suites generated by \texttt{gpt-4o-mini} (92 out of
120, or 76.7\%) never actually exercised the function they were supposed to test.

There are two gaps in current measurement practices that hide this problem. First, hardly
any study controls for the structural complexity of the code under test: results are
reported in aggregate, across functions with mixed and often unreported cyclomatic
complexity (CC)~\cite{wang2025testeval,jain2025testgeneval,zhang2024testbench}, so no one
knows at what point LLM-generated tests start to fail (GAP-T).
Because of this gap, we kept cyclomatic complexity fixed as a controlled variable, sampled
only the 5--10 band, and measured the CC--quality correlation directly (RQ3). Second,
branch coverage and mutation score are rarely combined in a single statistical test on a
controlled sample~\cite{sapozhnikov2024empirical,ouedraogo2025prompt,silva2025comparative}
(GAP-M), even though coverage alone can overstate quality --- in our own
baseline data, one \texttt{jsoup} function reaches 60\% branch coverage while killing 0\%
of mutants (Section~\ref{sec:discussion}).

Our design choices were shaped as much by lessons learned during implementation as by
insights from the literature. The original proposal named Defects4J and CodeXGLUE as data
sources; during preparation we found that CodeXGLUE is a code intelligence benchmark, not
a corpus of executable functions, so we mined the dataset ourselves: 120 functions
(60 Java, 60 Python) at pinned commits, all with a cyclomatic complexity between 5 and 10,
as measured by Lizard. A pilot run then showed that functions extracted from their modules
often cannot even resolve the names they reference --- in fact, 10 out of our 12 Python
pilot measurements were invalid for extraction-related reasons --- so in the final
experiment
every generated test runs inside its original project. This measurement approach is
important: measured this way, a score can no longer be reduced by artefacts from
extraction, only by the quality of the generated tests themselves. We compare
\texttt{gpt-4o-mini} (using one-shot prompting, temperature $0$) against three automated
baselines: EvoSuite and Randoop for Java (60 seconds per class), and Pynguin for Python
(90 seconds per module), scoring branch coverage and mutation score, using non-parametric
statistical tests with $\alpha = 0.05$.

We ask: \\
\textbf{RQ1}: Asked in a single one-shot prompt to test one function of medium cyclomatic
complexity ($5 \le CC \le 10$) drawn from a real open-source project, does GPT-4o-mini
produce a suite whose median branch coverage over that function's own lines reaches the
80\% industry target? \\
\textbf{RQ2}: Does the same suite reach a median mutation score of 60\% (RQ2-A), and does
it, matched function by function, score above the established automated generators ---
EvoSuite and Randoop for Java, Pynguin for Python (RQ2-B)? \\
\textbf{RQ3}: Within the narrow 5--10 band alone, does cyclomatic complexity correlate
negatively with the coverage the model achieves --- does quality degrade measurably as
complexity rises even inside one band?

These two thresholds are not arbitrary: 80\% is the branch coverage target most commonly
quoted as the industry standard, while 60\% is set intentionally above the roughly 34\%
mutation scores reported for LLM-generated suites in the studies reviewed in
Section~\ref{sec:related}; both were fixed in the approved proposal before any data were
collected.

To summarize, this paper contributes:
\begin{enumerate}
  \item A reproducible benchmark of 120 medium-complexity functions ($5 \le CC \le 10$) mined
    from ten pinned open-source projects, with full provenance for every function.
  \item A measurement pipeline that scores tests \emph{per function} inside the original
    projects and enforces a green-on-original check before any mutation analysis.
  \item An empirical comparison of GPT-4o-mini against search-based and random baselines on a
    complexity-matched sample, reporting branch coverage and mutation score together with
    effect sizes.
  \item A catalogue of the measurement pitfalls we encountered --- such as EvoSuite returning
    exit code 0 without generating a single test, or JaCoCo silently reporting 0\% under
    EvoSuite's instrumenting class loader --- which would have corrupted our results if left
    unnoticed.
\end{enumerate}

Section~\ref{sec:related} reviews prior work; Section~\ref{sec:method} details the
dataset, generation, and measurement process; Section~\ref{sec:results} answers each
research question; Section~\ref{sec:discussion} interprets our findings, including two
exploratory follow-ups; Section~\ref{sec:threats} discusses threats to validity;
Section~\ref{sec:conclusion} summarises our work.
